% \iffalse meta-comment
%
% Copyright (C) 2017- by Yanshuo Chu <yanshuoc@gmail.com>
%
% This file may be distributed and/or modified under the
% conditions of the LaTeX Project Public License, either version 1.3a
% of this license or (at your option) any later version.
% The latest version of this license is in:
%
% http://www.latex-project.org/lppl.txt
%
% and version 1.3a or later is part of all distributions of LaTeX
% version 2004/10/01 or later.
%
% \fi
%
% \section{版面与元素格式}
%
% \changes{v3.2a}{2026/09/20}{新版面的取值从 \file{hit-defaults.dtx} 拆出：页面设置从
%   “本科毕业论文书写范例（理工类）”解出，字号行距间距照写作指南 3.2 节，硕博与本科
%   共用；网格字距 249 缇、行距 391 缇（范例正文各章都是 391，383 只在拷贝进来的书写
%   规范那一章）；章标题与条标题 1.25 倍（指南四级都写 1.25）；标题与正文的默认值分
%   倍数版与毫米版两套由 \option{bottom} 挑；\option{foot-baseline} 1.52\,bp；人文社科
%   类正文固定值 23 磅不贴网格；深圳本科的开题与中期正文照终稿排、首页六段中英各一套}
%
% \option{layout} 与 \option{styles} 两族的取值。只有终稿这一档：报告还没有
% 新版面，它的版心写在 \file{src/flow/hit-flow-geometry.dtx} 里。
%
%    \begin{macrocode}
%<*hithesis-cfg>
%    \end{macrocode}
%
%    \begin{macrocode}
\hit@presetup@scope{body=final-body}
%    \end{macrocode}
%
% \subsection{本科的版面}
%
% 下面这一份本科与硕博共用。硕博不另解范例：两边的页面设置量出来一样，
% 字号行距也照同一份写作指南，只有图题表题的间距有出入，那一处等接
% \option{caption-figure} 与 \option{caption-table} 时再补。
%
% 页面设置来自学校发的
% “本科毕业论文书写范例（理工类）”，解法是拆开 \file{.docx} 读
% \texttt{w:sectPr} 与各个 \texttt{w:pPr}，再拿导出的 PDF 逐条基线核对，
% 量法记在 \file{src/docxlike/docxlike-page.dtx}；字号、行距、段前段后照
% “本科毕业论文（设计）写作指南（理工类）”3.2 节的原话。
% 威海本科没有对照文档，跟着这一份走。
%
% 条件写 \option{layout}\texttt{=!v2-layout}，也就是“没要旧版面”：本科写了
% \option{v2-layout}\texttt{=\{bachelor=true\}} 就回到 v3.1e 那套硬编码尺寸，
% 这一整段都不发，见 \file{src/flow/hit-flow-geometry.dtx}。
%
% \emph{页面设置。} 缇是 Word 的内部单位，一缇是 $1/20$\,bp：上边距 2155 缇、
% 左下右各 1701 缇、页眉距边界 1701 缇、页脚距边界 1304 缇、文档网格行跨度
% 383 缇。
%
% \emph{网格那两个数不是摆设。} Word 里“一行”“一字”这两个单位取的就是该节
% 文档网格的行跨度与字跨度：段前段后写“1 行”，存盘时折成
% “系数乘行跨度”向下取整到缇；首行缩进写“2 字”，折成 2 倍字跨度。
% 范例里每个章标题的段前都精确等于它所在那一节的 \texttt{w:linePitch}——
% 391 缇那一节的章标题写 391，383 缇那两节写 383，一个不差。所以
% \option{grid} 里行与字的 \option{pitch} 改一个数，标题间距与首行缩进
% 跟着全变，配置里那些“1 行”“2 字”一个字都不用动。
%
% 范例这一份文档十二节的网格并不统一：行跨度 391 缇八节、383 缇两节、
% 395 缇两节；字跨度 \texttt{w:charSpace} 1861 九节、2752 两节。取 391 与 1861：
% 正文各章都在 391 这一档，383 只在夹进来的书写规范那一章（那一节的章标题
% 段前 383 缇、段后 $0.8 \times 383 = 306.4$ 截成 306 缇，与它自己的网格对得上，
% 是拷贝进来时带着的），博士范例也是正文 391、拷贝进来的两章 383。
%
% 字距由 \texttt{w:charSpace} 定，那个数是字距超出字号的部分，单位是
% $1/4096$\,bp：1861 折成小四的 12\,bp 加 $1861/4096 = 0.4543$ 得 12.4543\,bp，
% 一行 34 字，首行缩进 2 字 24.91\,bp。另外两节写 2752（12.6719\,bp、一行
% 33 字、缩进 507 缇），结论那一章还写过 37928，那是“每行 20 字”的排版示范。
%
% \option{head} 是 \cs{headheight}，也就是页眉块的高度。这个量直接压在页眉
% 基线上：\hologo{LaTeX} 把页眉盒顶放在页眉距边界处，基线落在页眉距边界加
% \option{head} 再减去基线以下那一段的位置。基线以下是 7.4827\,bp——小五那
% 一行的深度 2.5326，加边框间距 1\,bp，加那条粗细双线的 2.25\,bp
% （\texttt{w:sz="18"} 是八分之一磅），加 \cs{headrule} 前后的零头。
% Word 那边页眉基线在页眉距边界下方一个 ascent 处，即 $1.0156 \times 9 =
% 9.1404$\,bp，所以 \option{head} 取 $9.1404 + 7.4827 = 16.6231$\,bp，量出来
% 页眉基线 94.19\,bp。范例是 94.08，它的坐标量化到 300\,dpi，一格 0.24\,bp，
% 读数本身就有半格的出入。
% 写作指南 3.4 节说双线“宽 0.8\,mm”，与 \cs{headrule} 里那条 2.276208\,pt
% 对得上。
% \option{foot-baseline} 的 1.52\,bp 是量出来的：页码基线在页脚线上方
% 1.52\,bp 处。范例的页码不是页脚段落，是个文本框
% （\texttt{<w:framePr w:hAnchor="margin" w:xAlign="center"/>}），位置由文本框
% 自己定，推不出闭式，只能量。
%
% 这个值越大页码越靠上：\file{src/docxlike/docxlike-page.dtx} 里
% \option{footskip} $=$ \option{margin-bottom} $-$ \option{footer} $-$
% \option{foot-baseline}，页脚基线在版心底下方 \option{footskip} 处。
% 先前取 2\,bp，改 2.48 是把它往上挪了两个像素，方向反了；现在退回来再往下
% 两个像素，净往下 0.96\,bp。
%
% \emph{标题。} 写作指南给了两份，一份按行数与多倍值（3.2 节），一份按毫米
% （3.8 节）。两份是并列的，不是一份加另一份的余量，所以这里也存成两套完整的
% 表，由 \option{layout}/\option{bottom} 挑。
%
% \option{bottom}\texttt{=ragged} 走倍数版，也就是 3.2 节那一份，逐条写死了，
% 照抄：
% \begin{longtblr}[caption={写作指南 3.2 节的标题版面（倍数版）}, label={tab:impl-layout-ragged}]
%   {colspec={l l l l}}
% \toprule
% & 字体字号 & 行距 & 段前段后 \\ \midrule
% 章标题 & 黑体小二 & 多倍 1.25 & 1 行 / 0.8 行 \\
% 节标题 & 黑体小三 & 多倍 1.25 & 0.5 行 / 0.5 行 \\
% 条标题 & 黑体四号 & 多倍 1.25 & 0.5 行 / 0.5 行 \\
% 款项标题 & 黑体小四 & 多倍 1.25 & 0 / 0 \\
% \bottomrule
% \end{longtblr}
% 四级都不对齐网格，章标题行距 $1.25 \times 1.296875 \times 18 = 29.18$\,bp。
% 范例的 heading 1 与 heading 3 样式写的是 1.2 倍（\texttt{w:line="288"}），
% 与指南原话不符，不跟；范例按指南改齐后让 Word 排的靶子（\texttt{make
% pages-doctor}）章到节 41.2、目录与第 1 章首页逐行对上。系数的来路见
% \file{src/docxlike/docxlike-format.dtx}。这一档全是确切值，没有可拉的东西，
% 所以版心底只能取 \texttt{ragged}。
%
% \option{bottom}\texttt{=flush} 走毫米版，也就是 3.8 节那一份。那里给的是
% “上、下行距”，指标题与前后两块之间那条看得见的空白，不是标题内部的行距：
% \begin{longtblr}[caption={写作指南 3.8 节的标题版面（毫米版）}, label={tab:impl-layout-flush}]
%   {colspec={l l l}}
% \toprule
% & 上、下行距 & 行距 \\ \midrule
% 章标题 & 10\,mm & 多倍 1.2\,--\,1.25 \\
% 节标题 & 7\,--\,8\,mm & 多倍 1.2\,--\,1.25 \\
% 条标题 & 6\,--\,7\,mm & 多倍 1.2\,--\,1.25 \\
% 款项标题 & 3\,--\,4\,mm & 多倍 1.2\,--\,1.25 \\
% 正文 & 3\,--\,4\,mm & —— \\
% \bottomrule
% \end{longtblr}
% 区间本身就是伸缩范围：小端做自然值，大端减小端做伸量，页面要撑到版心底
% 就往上拉。章标题原文只给一个数，没有区间，所以那一级的上下白是刚性的。
%
% 内部行距 3.8 节没给，四级一律沿用 3.2 节那两个值合成的区间 1.2\,--\,1.25，
% 不钉死其中一个。款项标题的 3\,--\,4\,mm 落在上下白上：它的段前段后本来就是
% 零，与前后正文之间那条白就是普通行距那一条。正文的 3\,--\,4\,mm 是
% \option{line-white}，量的也是行与行之间那条白。
%
% \emph{正文的倍数版与标题一样，多倍 1.25 加不对齐网格}，行距
% $1.25 \times 1.296875 \times 12 = 19.453$\,bp，首行基线落在版心顶下方
% $1.0078 \times 12 = 12.09$\,bp。本科与研究生这一条不同，别拿研究生
% 那边的范例来对：研究生是单倍加对齐网格，行距一格 19.55\,bp、首行落点
% 14.18\,bp，两个数只差半个百分点，一页三十三行下来差 3\,bp，看着像整体
% 偏了一点。
%
% \option{snap-to-grid} 关掉时，横向的字也不贴格子。Word 的这个开关管的是
% “对齐到网格”整件事，不只是行：本科的正文行是拉满版心的（范例满行末端
% 落在版心右边界 510.17\,bp），研究生那边贴格子，一行 34 个格占 423.45\,bp，
% 右边空着 1.73\,bp。这段余量由 \cs{g\_docxlike\_page\_slack\_dim} 算出来，
% 只在对齐网格时才发给 \cs{rightskip}，见 \file{src/docxlike/docxlike-format.dtx}。
%
% 范例有几处没照这一份排，都不取：前置部分两节的行跨度是 395 缇，摘要与目录
% 那几节的字跨度是 2752；结论那一章是“每行 20 字”的排版示范，字距 21.26\,bp，
% 正文也改成了对齐网格；致谢那一段也对齐网格。
%    \begin{macrocode}
\ExplSyntaxOff
\hit@presetup[layout=!v2-layout]{
  layout = {
    bottom = auto,
    margin = { top = 107.75bp, bottom = 85.05bp,
               left = 85.05bp, right = 85.05bp },
    header = { from-edge = 85.05bp,
               rule = { width = 2.25bp, from-text = 1bp } },
    footer = { from-edge = 65.2bp,
               shift = { vertical = -1.52bp, horizontal = 0.24bp } },
    grid = { lines = { pitch = 19.55bp },
             chars = { pitch = 12.4543bp, latin = true } },
    %% 按部件覆盖。范例逐节量下来：封面与内封的 docGrid 行跨度是 395 缇，
    %% 摘要、目录与正文 391。正文那一档就是上面的 19.55bp，不必重写。
    cover       = { grid = { lines = { pitch = 19.75bp } } },
    titlepage   = { grid = { lines = { pitch = 19.75bp } } },
  },
}
%% 右开页。规范没写要求右开，两份范例的分节符也全是 nextPage；博士封面那一串
%% 跳到右手页是双面装订的惯例，只有这一处。写 layout={openright=false} 即全篇
%% 不跳，交图书馆的电子版要的就是它。
\hit@presetup[degree-level=doctor]{
  layout = { cover = { openright-default = true } },
}
%% 本科声明页的两个小节标题带段前段后。范例 docx 段 662/668 写的是
%% w:before/after=195 缇，即网格行距 391 的一半，也就是 0.5 行；研究生那一份
%% （段 710/718）没有，间隔全由空段给。
\hit@presetup[degree-level=bachelor]{
  styles = {
    statement-heading     = { paragraph = { spacing = { before = 0.5 line, after = 0.5 line } } },
    authorization-heading = { paragraph = { spacing = { before = 0.5 line, after = 0.5 line } } },
  },
}
\hit@presetup[layout=!v2-layout]{
  styles = {
    Header = {
      font = { asian-text-font = SimSun, size = xiaowu, snap-to-grid = false },
      paragraph = { spacing = { line-spacing = multiple, at = 0 } },
    },
    Footer = {
      font = { size = xiaowu, snap-to-grid = false },
    },
    Normal = {
      font = { asian-text-font = SimSun, size = xiaosi, snap-to-grid = true },
      paragraph = { snap-to-grid = auto, indentation = { special = first-line, by = 2 chars } },
    },
    Normal/if-raggedbottom = { paragraph = { spacing = { line-spacing = multiple, at = 1.25 } } },
    Normal/if-flushbottom  = { paragraph = { spacing = { line-white = 3mm 4mm } } },
    Heading1 = {
      font = { asian-text-font = SimHei, size = xiaoer, snap-to-grid = false, character-spacing = { spacing = expanded, by = 0.0543bp } },
      paragraph = { alignment = centered },
    },
    Heading1/if-raggedbottom = { paragraph = { spacing = { line-spacing = multiple, at = 1.25, before = 1 line, after = 0.8 line } } },
    Heading1/if-flushbottom  = { paragraph = { spacing = { line-range = {1.2 1.25}, white-before = 10mm, white-after = 10mm } } },
    Heading2 = {
      font = { asian-text-font = SimHei, size = xiaosan, snap-to-grid = false, character-spacing = { spacing = expanded, by = 0.4543bp } },
    },
    Heading2/if-raggedbottom = { paragraph = { spacing = { line-spacing = multiple, at = 1.25, before = 0.5 line, after = 0.5 line } } },
    Heading2/if-flushbottom  = { paragraph = { spacing = { line-range = {1.2 1.25}, white-before = 7mm 8mm, white-after = 7mm 8mm } } },
    Heading3 = {
      font = { asian-text-font = SimHei, size = sihao, snap-to-grid = false, character-spacing = { spacing = expanded, by = 0.4543bp } },
    },
    Heading3/if-raggedbottom = { paragraph = { spacing = { line-spacing = multiple, at = 1.25, before = 0.5 line, after = 0.5 line } } },
    Heading3/if-flushbottom  = { paragraph = { spacing = { line-range = {1.2 1.25}, white-before = 6mm 7mm, white-after = 6mm 7mm } } },
    Heading4 = {
      font = { asian-text-font = SimHei, size = xiaosi, snap-to-grid = false, character-spacing = { spacing = expanded, by = 0.4543bp } },
    },
    Heading4/if-raggedbottom = { paragraph = { spacing = { line-spacing = multiple, at = 1.25, before = 0bp, after = 0bp } } },
    Heading4/if-flushbottom  = { paragraph = { spacing = { line-range = {1.2 1.25}, white-before = 3mm 4mm, white-after = 3mm 4mm } } },
    abstract-en = {
      font = { latin-line = true, latin-text-font = Times New Roman, size = xiaosi, snap-to-grid = false },
      paragraph = { snap-to-grid = false, spacing = { line-spacing = multiple, at = 1.25 } },
    },
    %% 题注。字号由写作指南 2.12／2.13 写死（“表题…用宋体5号字”「图题…用
    %% 宋体5号字」）；行距规范一个字没提，取 1.25，与正文同一档。
    %%
    %% 范例里题注那几段全是排版的人一段一段手敲的直接格式，自己就打架：同一份
    %% 文件里图1-1 第一处无 spacing 即单倍、第二处 w:line=280 exact 即固定 14pt，
    %% 表6-1 是 1.2 倍、表6-2 是 1.25 倍。styles.xml 里根本没有“题注”这个样式，
    %% 没有一处文档级的规矩可依。判据的次序是样式定义、指南原话、段落级直接
    %% 格式；题注这几段全在第三档，1.25 是那几个值里唯一与正文一致的，范例
    %% 自己也用过。
    %%
    %% 题与图表之间不给间距（\abovecaptionskip 与 \belowcaptionskip 都是零）。
    %% 范例里图题紧跟着图，图片下沿量不到，没有干净的实测，先取零。
    %%
    %% 横向不写 snap-to-grid，跟着正文走，也就是照样贴字符网格。字符网格的格
    %% 宽是 12.4543bp，五号字身 10.5bp，一格里的增量 1.9543bp 比正文那一档
    %% （小四，0.4543bp）大，所以同一行题注比接进版面之前宽一截。这是网格模型
    %% 本身的结果，不是这里另加的间距：半角字在网格里吃半格，两个半角正好一个
    %% 全角，规范说的“图名在图序之后空 2 个半角字符”在这个模型里才是一个字宽。
    caption-figure = { font = { asian-text-font = SimSun, size = wuhao }, paragraph = { spacing = { line-spacing = multiple, at = 1.25 } } },
    caption-table  = { font = { asian-text-font = SimSun, size = wuhao }, paragraph = { spacing = { line-spacing = multiple, at = 1.25 } } },
    caption-sub    = { font = { asian-text-font = SimSun, size = wuhao }, paragraph = { spacing = { line-spacing = multiple, at = 1.25 } } },
    %% 图注。规范把它归在“图中文字”那一句里，与题注同一档字号行距。
    caption-note   = { font = { asian-text-font = SimSun, size = wuhao }, paragraph = { spacing = { line-spacing = multiple, at = 1.25 } } },
    %% 表格里的段。规范“表中文字用宋体 5 号字”；行高一个网格行、行盒在行里
    %% 居中，是本科范例第 12 页那张表量出来的（表头下 387.12、底线 521.40、
    %% 七行，一行 19.18，那一节 linePitch=383 即 19.15），也就是单倍贴格。
    table-cell     = {
      font = { asian-text-font = SimSun, size = wuhao },
      paragraph = { snap-to-grid = true, spacing = { line-spacing = multiple, at = 1 } },
    },
    titlepage-mark = {
      font = { asian-text-font = SimSun, size = xiaosi, snap-to-grid = false },
      paragraph = { snap-to-grid = false },
    },
    titlepage-label = {
      font = { asian-text-font = SimSun, size = xiaoer, font-style = bold, snap-to-grid = false },
      paragraph = { alignment = centered, snap-to-grid = false },
    },
    titlepage-title = {
      font = { asian-text-font = SimHei, size = erhao, snap-to-grid = false },
      paragraph = { alignment = centered, snap-to-grid = true },
    },
    titlepage-field = {
      font = { asian-text-font = SimSun, size = sihao, snap-to-grid = false },
      paragraph = { alignment = centered, snap-to-grid = true, spacing = { line-spacing = multiple, at = 1.5 } },
    },
    %% 英文内封。取值来自研究生范例 docx 段 87-115：分类号两行单倍；
    %% “Dissertation for the …”小二；题目 w:line=300 即 1.25 倍、不贴网格；
    %% 字段是 7 行 2 列的表，w:line="440" w:lineRule="exact" 即固定值 22pt。
    titlepage-en-mark = {
      font = { latin-line = true, latin-text-font = Times New Roman, size = xiaosi, snap-to-grid = false },
      paragraph = { snap-to-grid = true, spacing = { line-spacing = multiple, at = 1 } },
    },
    titlepage-en-degree = {
      font = { latin-line = true, latin-text-font = Times New Roman, size = xiaoer, snap-to-grid = false },
      paragraph = { alignment = centered, snap-to-grid = true, spacing = { line-spacing = multiple, at = 1 } },
    },
    titlepage-en-title = {
      font = { latin-line = true, latin-text-font = Times New Roman, size = xiaoer, font-style = bold, all-caps = true, snap-to-grid = false },
      paragraph = { alignment = centered, snap-to-grid = false, spacing = { line-spacing = multiple, at = 1.25 } },
    },
    titlepage-en-field = {
      font = { latin-line = true, latin-text-font = Times New Roman, size = sihao, snap-to-grid = false },
      paragraph = { snap-to-grid = false, spacing = { line-spacing = exactly, at = 22bp } },
    },
    %% 声明页里的两个小节标题。研究生范例 docx 段 710/718：sz30 即小三，黑体，
    %% 居中，line=300 即 1.25 倍，不贴网格，没有段前段后——间隔由空段给。
    %% 本科那一份的段前段后是 195 缇即 0.5 行，见下面按学位的预设。
    %%
    %% 两节的取值眼下一模一样，仍分成两个目标：相同是“现在恰好相同”，不是
    %% “必须相同”。名字与 term 族里那两条同名——一个说它写什么，一个说它长
    %% 什么样。
    statement-heading = {
      font = { asian-text-font = SimHei, size = xiaosan, snap-to-grid = false },
      paragraph = { alignment = centered, snap-to-grid = false, spacing = { line-spacing = multiple, at = 1.25 } },
    },
    authorization-heading = {
      font = { asian-text-font = SimHei, size = xiaosan, snap-to-grid = false },
      paragraph = { alignment = centered, snap-to-grid = false, spacing = { line-spacing = multiple, at = 1.25 } },
    },
    cover-label = {
      font = { asian-text-font = SimSun, size = xiaoyi, font-style = bold, snap-to-grid = false },
      paragraph = { alignment = centered, snap-to-grid = false, spacing = { line-spacing = multiple, at = 1 } },
    },
    cover-type = {
      font = { asian-text-font = SimSun, size = xiaoer, font-style = bold, snap-to-grid = false },
      paragraph = { alignment = centered, snap-to-grid = false, spacing = { line-spacing = multiple, at = 1 } },
    },
    cover-title = {
      font = { asian-text-font = SimHei, size = erhao, snap-to-grid = false },
      paragraph = { alignment = centered, snap-to-grid = true, spacing = { line-spacing = multiple, at = 1, before = 12bp } },
    },
    cover-subtitle = {
      font = { asian-text-font = SimHei, size = xiaoer, snap-to-grid = false },
      paragraph = { alignment = centered, snap-to-grid = false, spacing = { line-spacing = multiple, at = 1 } },
    },
    cover-title-en = {
      font = { latin-line = true, latin-text-font = Times New Roman, size = xiaoer, font-style = bold, all-caps = true, snap-to-grid = false },
      paragraph = { alignment = centered, snap-to-grid = false, spacing = { line-spacing = multiple, at = 1.25, before = 12bp } },
    },
    cover-author = {
      font = { asian-text-font = SimSun, size = xiaoer, font-style = bold, snap-to-grid = false },
      paragraph = { alignment = centered, snap-to-grid = false, spacing = { line-spacing = multiple, at = 1 } },
    },
    cover-institution = {
      font = { asian-text-font = KaiTi, size = xiaoer, font-style = bold, snap-to-grid = false },
      paragraph = { alignment = centered, snap-to-grid = false, spacing = { line-spacing = multiple, at = 1.35 } },
    },
    cover-date = {
      font = { asian-text-font = SimSun, size = xiaoer, font-style = bold, snap-to-grid = false },
      paragraph = { alignment = centered, snap-to-grid = false, spacing = { line-spacing = multiple, at = 1 } },
    },
  },
}
%% 目录条目四级。字号小四、一级黑体其余宋体，取自两份范例的 TOC1-TOC4 样式；
%% 左缩进 0／240／480／720 缇由 tocloft 那边给，这里只写算行高要用的三项。
%% 行距是范例的多倍 1.25，研究生那一档 1.2、人文社科类固定值 23bp，分别在下面
%% 两处覆盖。消费端是 \docxlike_format_stack_line:n，见 src/docxlike/docxlike-format.dtx。
%%
%% 这四条不挂 layout=!v2-layout。目录条目的贴排制两套版面都在走，旧版面那边
%% 也是照 Word 排的；挂上去旧版面就没有字号与行距可读了。
\hit@presetup{
  styles = {
    toc-chapter = {
      font = { asian-text-font = SimHei, size = xiaosi, snap-to-grid = false },
      paragraph = { snap-to-grid = false, spacing = { line-spacing = multiple, at = 1.25 } },
    },
    toc-section = {
      font = { asian-text-font = SimSun, size = xiaosi, snap-to-grid = false },
      paragraph = { snap-to-grid = false, spacing = { line-spacing = multiple, at = 1.25 } },
    },
    toc-subsection = {
      font = { asian-text-font = SimSun, size = xiaosi, snap-to-grid = false },
      paragraph = { snap-to-grid = false, spacing = { line-spacing = multiple, at = 1.25 } },
    },
    toc-subsubsection = {
      font = { asian-text-font = SimSun, size = xiaosi, snap-to-grid = false },
      paragraph = { snap-to-grid = false, spacing = { line-spacing = multiple, at = 1.25 } },
    },
  },
}
%% 英文学位论文的章标题全大写（“CHAPTER 1”“INTRODUCTION”“ABSTRACT”），是 Word 的
%% “全部大写字母”，编号跟着段落标记一起大写。报告的顶级标题不大写。
\hit@presetup[stage=final,lang=en,layout=!v2-layout]{
  styles = { Heading1 = { font = { all-caps = true } } },
}
%% 研究生范例的目录条目是多倍 1.2，本科那份 1.25，别的逐项相同：两份 docx 的
%% TOC1／TOC2／TOC3 样式除一个 w:lang 外逐字节一样，页面与文档网格、右制表位、
%% 四级缩进也一样。两边是从不同方向凑到各自那个数的——本科 TOC1 用样式的
%% line=300，TOC2／TOC3 在段上直接格式覆盖成 300；研究生 TOC1／TOC2 段上写
%% 288，TOC3 用样式的 288。
\hit@presetup[degree-level=!bachelor]{
  styles = {
    toc-chapter       = { paragraph = { spacing = { line-spacing = multiple, at = 1.2 } } },
    toc-section       = { paragraph = { spacing = { line-spacing = multiple, at = 1.2 } } },
    toc-subsection    = { paragraph = { spacing = { line-spacing = multiple, at = 1.2 } } },
    toc-subsubsection = { paragraph = { spacing = { line-spacing = multiple, at = 1.2 } } },
  },
}
%% 本科内封头一行（勾选框那一行）吃网格：范例 docx 里它没写 snapToGrid=0，
%% 首行基线落在版心顶下 14.25（小四上升部 12.18 加半格余量 2.09），不是 12.18；
%% 研究生那一份写了 snapToGrid=0，首行 119.8，两边不同档。写在总表之后，
%% 不然被总表里那条不贴格的盖掉。
\hit@presetup[degree-level=bachelor]{
  styles = { titlepage-mark = { paragraph = { snap-to-grid = true } } },
}
\ExplSyntaxOn
%    \end{macrocode}
%
% \subsection{人文社科类}
%
% 人文社科类的正文行距。指南写的是 1.25 倍，两份人文范例（博士 \texttt{s-hss}、
% 本科）的正文段却是 \texttt{line=460 exact} 加 \texttt{snapToGrid=0}，Word
% 导出逐行量步进一律 23.0；照范例。目录条目同为 23 磅固定值，四级各写一条，
% 与正文那条走同一个模型。
%
% 英文目录的四处字样与理工范例不同：“Abstract(In Chinese)”括号前不空、
% “Acknowledgments”没有 e、“letter”小写，都是人文博士范例英文目录里的写法。
%    \begin{macrocode}
\ExplSyntaxOff
\hit@presetup[category=hass]{
  styles = {
    Normal = { paragraph = { snap-to-grid = false } },
    Normal/if-raggedbottom = { paragraph = { spacing = { line-spacing = exactly, at = 23bp } } },
    toc-chapter       = { paragraph = { spacing = { line-spacing = exactly, at = 23bp } } },
    toc-section       = { paragraph = { spacing = { line-spacing = exactly, at = 23bp } } },
    toc-subsection    = { paragraph = { spacing = { line-spacing = exactly, at = 23bp } } },
    toc-subsubsection = { paragraph = { spacing = { line-spacing = exactly, at = 23bp } } },
  },
}
\ExplSyntaxOn
%    \end{macrocode}
%
%    \begin{macrocode}
%    \end{macrocode}
%
% \subsection{深圳本科那两份报告的差额}
%
% \emph{从终稿那一份起，只覆盖四处。} 学校发的《深圳本科毕业论文（设计）开题报告》
% 与《中期报告》两份原件，三节的 \texttt{w:pgMar} 与 \texttt{w:docGrid} 逐个数
% 与终稿相同，正文那一节正是 2154／1701 加 391／1861。差的只有下面四条。
% 这四个数取自 iota-hit 的 \texttt{src/config/layout.typ}，那边是拆开
% \file{.docx} 量出来的，本仓库没有原件，没有复核过。
%
% \emph{网格不启用。} 那一节的 \texttt{docGrid} 写着 391／1861，可没有
% \texttt{w:type}——量出来正文的汉字步进是 12.0，不是终稿的 12.4543；行距是
% 1.25 倍的 19.44，不是网格的 19.55。数照留，只说它不生效，那件事由
% \option{grid}/\option{enabled} 管。
%
% \emph{不启用还管着“一行”有多长。} 章节标题的段前段后写的是“1 行”“0.8 行”
% “0.5 行”，网格关着时折成正文的一个字号，即 12／9.6／6\,bp，与原件段落上缓存的
% 240／192／120 缇一个不差；按网格行距折会得到 19.15／15.32／9.58。
%
% \emph{页眉页脚都印。} 页眉距边界 2.4\,cm（\texttt{w:header="1361"}）加那一段
% 自己写的段前 12\,bp，横线 \texttt{thinThickSmallGap} \texttt{sz24} 即 3\,bp、
% 离字 4\,bp（\texttt{w:space="4"}）。
%    \begin{macrocode}
\hit@presetup@scope{body=final-body}
\ExplSyntaxOff
\hit@presetup[campus=shenzhen,degree-level=bachelor]{
  layout = {
    grid = { enabled = false ,
             lines = { pitch = 19.55bp } , chars = { pitch = 12bp } },
    header = { from-edge = 80.03bp,
               rule = { width = 3bp, from-text = 4bp } },
  },
}
\hit@presetup[campus=shenzhen,degree-level=bachelor]{
  styles = {
    Normal = { font = { snap-to-grid = false }, paragraph = { snap-to-grid = false } },
  },
}
\ExplSyntaxOn
%    \end{macrocode}
%
% \emph{深圳本科报告首页的段。} 那一页自己排自己的，见
% \file{src/pages/hit-report-cover-shenzhen-bachelor.dtx}；这里只给六段的版面。
% 数取自 iota-hit，那边是拆 2026 年 9 月版两份原件的 \texttt{.docx} 量出来的。
%
% 头一段是勾选行。字是小四不是小一：原件那一段的 \texttt{w:rPr} 写着
% \texttt{sz="48"}，可那是\emph{段落标记}的字号，两个框与四个字的 run 一个
% \texttt{sz} 都没写，跟 Normal 走。行盒 19\,bp 固定值，与字号无关。
%
% 第二段空着，中文版二号行距 1.25 倍，英文版与勾选行同一档。首页那一节的
% \texttt{docGrid} 写着 395／1861 却没有 \texttt{w:type}，整节的网格都不启用，
% 所以这几段一律不贴网格。
%
% 校名与文种中文版都是小初华文新魏加粗、行距 1.2 倍；英文版的文种改成宋体
% 小一加粗、段前段后各 3.9\,bp，另在两行下面各排一行 Times 加粗的英文，
% 校名 20\,bp 行距 1.25 倍、文种 16\,bp。
%
% 这一页自成一节，页边距与正文那一节不同，所以 \option{cover} 那一档部件差额
% 也写在这里。
%
% \emph{英文版还有三处与中文版不同。} 正文首行缩进三个字；章与节在原件里是
% Times Bold，中文版一级都不加粗；目录条目的行距 1.5 倍。三条都在下面写出来了。

%    \begin{macrocode}
\ExplSyntaxOff
\hit@presetup@scope{stage=!final}
%% 破折号排西文那一副。逐份数深圳那几份报告原件里破折号 run 上的 w:hint：
%% 中文中期那份 11 个一个 eastAsia 都没有，英文中期同样，Word 排出来就是 Times
%% 的那一条杠（中易宋体的两个“——”中间断开一截）。开题两份里没有破折号，
%% 跟同分支的中期；深圳硕博四份表单也没有，跟同校区的本科报告。
%% 学位论文那两份范例里 eastAsia 占多数，所以只给报告这一档。
\hit@presetup[campus=shenzhen]{ style = { em-dash = latin } }
%% 首页自成一节，两个阶段同一套（2026 年 9 月版两份逐个数相同）：
%% w:pgMar 1440/1800/1440/1800 即上下 72bp、左右 90bp；docGrid 395/1861 同样
%% 没有 w:type，网格不启用。页眉页脚那一节都空着（w:titlePg）。
\hit@presetup[campus=shenzhen,degree-level=bachelor]{
  layout = {
    cover = {
      margin = { top = 72bp, bottom = 72bp, left = 90bp, right = 90bp },
      grid = { lines = { pitch = 19.75bp } , chars = { pitch = 12bp } },
    },
  },
}
\hit@presetup[campus=shenzhen,degree-level=bachelor]{
  styles = {
    %% 勾选行顶格排，不居中：原件那一段的 w:jc 没写，Word 默认左对齐，
    %% 头一个方框的左缘就是版心左缘 90.0（逐字比 iota 的渲染件：方框宽
    %% 0.891em＝10.69，“毕”落在 100.69）。
    report-cover-form = {
      font = { asian-text-font = SimSun, size = xiaosi, snap-to-grid = false },
      paragraph = { snap-to-grid = false, spacing = { line-spacing = exactly, at = 19bp, before = 12bp } },
    },
    %% 空段的高按段落标记算，而段落标记取 w:rFonts 的 ascii 槽即西文那一副，
    %% 行高系数 1.15 不是中易的 1.296875。写 font-zh=none 就是这一档。
    report-cover-blank = {
      font = { latin-line = true, size = erhao, snap-to-grid = false },
      paragraph = { alignment = centered, snap-to-grid = false, spacing = { line-spacing = multiple, at = 1.25, before = 12bp } },
    },
    report-cover-university = {
      font = { asian-text-font = STXinwei, size = xiaochu, font-style = bold, snap-to-grid = false },
      paragraph = { alignment = centered, snap-to-grid = false, spacing = { line-spacing = multiple, at = 1.2 } },
    },
    report-cover-stage = {
      font = { asian-text-font = STXinwei, size = xiaochu, font-style = bold, snap-to-grid = false },
      paragraph = { alignment = centered, snap-to-grid = false, spacing = { line-spacing = multiple, at = 1.2 } },
    },
    %% 表里一格一段：段前 12bp、行距 1.5 倍。六格的行盒由标签与值里（英文版
    %% 标签带 16bp 的尾巴）高的那个定，所以字号写在这儿的是行盒那一档，
    %% 标签与值自己的字号在排版代码里发。
    report-cover-field = {
      font = { asian-text-font = SimSun, size = xiaosan, snap-to-grid = false },
      paragraph = { snap-to-grid = false, spacing = { line-spacing = multiple, at = 1.5, before = 12bp } },
    },
    report-cover-title = {
      font = { asian-text-font = SimHei, size = xiaoer, snap-to-grid = false },
      paragraph = { snap-to-grid = false, spacing = { line-spacing = multiple, at = 1.5, before = 12bp } },
    },
  },
}
\hit@presetup[campus=shenzhen,degree-level=bachelor,lang=en]{
  styles = {
    report-cover-blank = { font = { size = xiaoyi }, paragraph = { spacing = { line-spacing = exactly, at = 19bp } } },
    report-cover-university = { paragraph = { spacing = { before = 12bp } } },
    report-cover-university-en = {
      font = { latin-line = true, latin-text-font = Times New Roman, size = 20bp, font-style = bold, snap-to-grid = false },
      paragraph = { alignment = centered, snap-to-grid = false, spacing = { line-spacing = multiple, at = 1.25, before = 12bp } },
    },
    report-cover-stage = {
      font = { asian-text-font = SimSun, size = xiaoyi },
      paragraph = { spacing = { line-spacing = multiple, at = 1, before = 3.9bp, after = 3.9bp } },
    },
    %% 段前写零：Word 相邻两段取段后与段前的较大者，而 \cs{docxlike\_format\_par:nn}
    %% 是两段胶各发各的、会相加。上一段的段后正好也是 3.9，留它一个就对。
    report-cover-stage-en = {
      font = { latin-line = true, latin-text-font = Times New Roman, size = 16bp, font-style = bold, snap-to-grid = false },
      paragraph = { alignment = centered, snap-to-grid = false, spacing = { line-spacing = multiple, at = 1, before = 0bp, after = 3.9bp } },
    },
    %% 校名与文种之间还有一个空段，段落标记只有 8bp（w:sz="16"），
    %% 行距 1.25 倍、段前 12bp。中文版没有这一段。
    report-cover-blank-mid = {
      font = { latin-line = true, size = 8bp, snap-to-grid = false },
      paragraph = { alignment = centered, snap-to-grid = false, spacing = { line-spacing = multiple, at = 1.25, before = 12bp } },
    },
    %% 英文版标签带一截 16bp 的“/ English”，一行的行盒取这一行最大的那个 run，
    %% 六格的步进因此是 43.13 而不是 41.18（原件量得 43.2）。
    report-cover-field = { font = { size = 16bp } },
    %% 英文版正文首行缩进三个字（原件如此），中文版与终稿一样两个。
    %% 行盒按西文算，理由同下面的章与节：英文版的正文段整段没有汉字，
    %% 上升部是 Times 的而不是宋体的，小四 12bp 下差 0.90。
    Normal = { font = { latin-line = true }, paragraph = { indentation = { special = first-line, by = 3 chars } } },
    %% 两处没写，写了也不动，记在这里省得再试一遍：
    %%
    %% 页眉的行盒高深是钉死的常量（hit-final-pagestyle.dtx 的
    %% \c__hit_header_text_ht_dim），与 format 无关。原件里英文页眉该不该比
    %% 中文的高 0.65，手上没有那两份英文原件可量。
    %%
    %% 题注这一档报告根本读不到：报告的 \@makecaption 是 ccaption 的，
    %% 只有终稿那一份 \@makecaption 才调 \hit@caption@font: 去读
    %% caption-figure／caption-table。正文照终稿排的那一档题注也该照终稿排，
    %% 那是另一批的事。按模型算，英文版的题注上升部该是 9.874（五号 Times），
    %% 与 iota-hit 渲染件反解出来的一致。
    %% 英文版的题注上下间距。原件里图题走 caption 样式、段上写着
    %% w:before="0" w:after="60"（段后 3bp）；表题走“题注-tab”样式、段上写着
    %% w:beforeLines="20“（0.2 行 ＝ 2.4bp），样式里 w:after=”0"。
    caption-figure = { paragraph = { spacing = { before = 0bp, after = 3bp } } },
    caption-table  = { paragraph = { spacing = { before = 2.4bp, after = 0bp } } },
    %% 英文版目录条目 1.5 倍。中期原件的 TOC1／TOC2／TOC3 都写着 w:line="360"，
    %% 开题的 TOC1／TOC2 也是 360、TOC3 却是 300；三级条目一共三行，一份目录里
    %% 一二级 1.5、三级 1.25 说不出道理，当漏改论，三级照 1.5。中文版那两份
    %% 三级都写 300，与终稿本科那一档同数，不必另列。
    %% 只有深圳本科这两份验过，键名里没写 -shenzhen 是因为本部没发过本科的
    %% 英文报告；本部哪天发了，这一条要么加校区、要么按那一份重来。
    toc-chapter       = { paragraph = { spacing = { line-spacing = multiple, at = 1.5 } } },
    toc-section       = { paragraph = { spacing = { line-spacing = multiple, at = 1.5 } } },
    toc-subsection    = { paragraph = { spacing = { line-spacing = multiple, at = 1.5 } } },
    toc-subsubsection = { paragraph = { spacing = { line-spacing = multiple, at = 1.5 } } },
    %% 英文版的章与节在原件里是 Times Bold，中文版一级都不加粗。
    %% 只加粗西文那一槽：黑体没有粗体面，整段加粗会让夹在英文标题里的汉字
    %% 走描边合成，见本类 hit-docxlike.dtx 的 \cs{\_\_hit\_docxlike\_latin\_bold:}。
    %%
    %% 行盒也按西文算。Word 的行盒由这一行内容的字体定，英文版的标题整行没有
    %% 汉字，上升部是 Times 的 0.9336 而不是宋体的 1.0078：小二 18bp 下差 1.36
    %% （iota-hit 的渲染件里章标题基线 136.62，中文版同一处 137.98）。
    %% 写 latin-line 只改行盒，asian-text-font 留着，标题里万一夹了汉字还是黑体。
    %%
    %% iota-hit 那边的判据是“这一段有没有汉字”，逐条标题现判；这边按语种给，
    %% 英文档里排一条中文标题就判错。要逐条判得先有个拿得到标题正文的钩子，
    %% ctex 的 format 收不到标题，另记。
    Heading1 = { font = { latin-bold = true, latin-line = true } },
    Heading2 = { font = { latin-bold = true, latin-line = true } },
  },
}
%% 章标题与终稿差两条。行距 1.25：那两份的说明页写着「章标题　黑体小 2 号字，
%% 多倍行距值 1.25」，原件的段落也是 w:line="300" auto；终稿那一档是 1.2（范例如此），
%% 两边差在行深上，18bp 的行深 10.95 对 9.79，章到节因此差 1.17。
%% 字距 0：终稿范例的章标题 run 上写着 w:spacing="-8"，报告那几份一个都没写。
%% 行距要写进两个档：终稿那边写的是 chapter/if-raggedbottom，分档的取值压过
%% 不分档的，光写 chapter 盖不住。
%% 题注上下的间距。中文那两份原件里图段、图题、表题、表格全是光
%% w:line="300" auto snapToGrid=0，段上一个 w:before／w:after 都没写，
%% 四个空当都是零；英文那一份另有取值，见上面 lang=en 那一块。
\hit@presetup[campus=shenzhen,degree-level=bachelor]{
  styles = {
    caption-figure = { paragraph = { spacing = { before = 0bp, after = 0bp } } },
    caption-table  = { paragraph = { spacing = { before = 0bp, after = 0bp } } },
  },
}
\hit@presetup[campus=shenzhen,degree-level=bachelor]{
  styles = {
    Heading1 = { font = { character-spacing = { spacing = normal } } },
    Heading1/if-raggedbottom = { paragraph = { spacing = { line-spacing = multiple, at = 1.25 } } },
    Heading1/if-flushbottom  = { paragraph = { spacing = { line-spacing = multiple, at = 1.25 } } },
  },
}
%% 表格。中期报告原件那张表（表 2-1）每一格的段都是 line=300 auto 加
%% snapToGrid=0，字号没另写就是正文的小四，实测数据行步进 19.44，正是小四
%% 1.25 倍的行盒；不是学位论文那档的五号贴格。图表前后原件一个 w:before／
%% w:after 都没写，空一行是作者自己敲的空段（\hitenter），浮动体间距取零。
\hit@presetup[campus=shenzhen,degree-level=bachelor]{
  styles = {
    table-cell = {
      font = { size = xiaosi, snap-to-grid = false },
      paragraph = { snap-to-grid = false, spacing = { line-spacing = multiple, at = 1.25 } },
    },
  },
  style = { float-sep = 0bp },
}
\ExplSyntaxOn
%    \end{macrocode}
%
%    \begin{macrocode}
\hit@presetup@scope{}
%    \end{macrocode}
%
%    \begin{macrocode}
%</hithesis-cfg>
%    \end{macrocode}
%
% \Finale
